\scanpage{88}
% archive sheet number: 87
% block text | lines 1-2 cancelled block
\deleted{aussi, et par suite aussi la première flèche \deleted{est} horizontale.}

% block text | lines 3-4 cancelled block
\deleted{\deleted{\unsure{Supposons}} Prouvons donc que $\varphi$ est injective. Prouvons par récurrence sur $n\geq1$ que \deleted{\unsure{si}}}
% block display | lines 5-6 cancelled block
\[
\begin{aligned}
&\deleted{\text{si}\qquad f\sim g\quad\bmod\unsure{\mathfrak r|\mathfrak m|\mathfrak b}^n A}\\
&\deleted{\text{alors}\qquad f\sim g\quad\bmod\mathfrak r^{n+1}A}
\end{aligned}
\]
% block text | line 7 cancelled block
\deleted{Pour ceci considérons la suite exacte des $\Gamma$-groupes}
% block diagram | line 8 cancelled block
\diagram{\begin{tikzcd}[ampersand replacement=\&,column sep=small]
\deleted{e}\arrow[r]\&\deleted{N_n}\arrow[r]\&\deleted{(A/\mathfrak r^{n+1}A)^*}\arrow[r]\&\deleted{(A/\mathfrak r^n A)^*}\arrow[r]\&\deleted{0}
\end{tikzcd}}
% block text | lines 9-12 cancelled block
\deleted{Si $x,y\in A^*\deleted{\unsure{/\mathfrak r^{n+1}A}}$ leurs images dans $A/\mathfrak r^n A$ sont égales si et seulement si $x-y\in\mathfrak r^n A$ i.e. $xy^{-1}-1\in\mathfrak r^n A$ \deleted{et de $=$} i.e. $xy^{-1}\in1+\mathfrak r^n A$, et de $=$ pour l'\unsure{égalité|identité} \deleted{\unsure{suivant}} mod $\mathfrak r^{n+1}A$.}

% block text | lines 13-14 cancelled block
\deleted{Donc $N_n$ s'identifie au groupe additif $\mathfrak r^n A/\mathfrak r^{n+1}A$}
% block text | lines 14-15 restart
Je veux prouver que si $f,g$ sont tels que
% block display | line 16
\[
g(\unsure{\gamma|\sigma})=u^{-1}f(\gamma)\gamma(u)\qquad\gamma\in\Gamma_i\tag*{(*)}
\]
% block text | lines 17-18
alors $f\sim g$, \inserted{\unsure{i.e.} \unsure{$-$}} Nous pouvons d'abord faire opérer $\Gamma$ différemment sur $A$, par
% block display | line 19
\[
\gamma_f(x)=f(\gamma)\gamma\cdot x f(\gamma)^{-1}
\]
% block text | line 20 left
On a donc
% block display | line 20
\[
\gamma_f(\lambda x)=f(\gamma)\gamma(\lambda)\gamma(x)f(\gamma)^{-1}=\gamma(\lambda)\gamma_f(x)
\]
% block text | lines 21-22
donc $A_f$ est un $(O_f,\Gamma)$-module. \deleted{\unsure{Donc|Dans}} Cela étant, les équations (*) \unsure{s'écrit|s'écrivent}
% block display | line 23
\[
g(\gamma)f(\gamma)^{-1}=u^{-1}\gamma_f(u)\qquad\gamma\in\Gamma_i
\]
% block text | line 24
et on cherche $v$ tel qu'on ait
% block display | line 25
\[
g(\gamma)f(\gamma)^{-1}=v^{-1}\gamma_f(v)\qquad\text{pour }\text{\uuline{$\unsure{\mathrm{tt}|\forall}$}}\ \gamma\in\Gamma.
\]
